\chapter[Execução do Experimento]{Execução do Experimento}

\section{Visão geral da execução}

Este capítulo descreve a execução do experimento conduzido neste trabalho, detalhando as etapas realizadas para operacionalizar 
o planejamento apresentado no capítulo anterior. A execução envolveu a preparação da base experimental, 
a organização dos instrumentos de coleta, a configuração do ambiente utilizado, a aplicação dos trechos e prompts ao SonarQube e modelos de linguagem 
avaliados respectivamente.

A partir das amostras selecionadas do dataset MLCQ e dos vereditos de referência definidos no ground truth, foram 
construídas as entradas submetidas às LLMs. Cada entrada foi associada a um identificador de prompt, ao tipo de elemento 
analisado, ao code smell avaliado e à respectiva versão de prompt. Em seguida, as respostas geradas pelos modelos foram coletadas, 
padronizadas e integradas à base de referência, permitindo a posterior comparação entre os vereditos produzidos pelas LLMs e os vereditos esperados.

Também são descritos neste capítulo os procedimentos adotados para validação dos dados coletados, incluindo a verificação de respostas ausentes, 
registros duplicados, inconsistências de identificação e padronização dos rótulos extraídos. Dessa forma, a execução do experimento resultou 
em uma base analítica estruturada, utilizada nas etapas posteriores de análise descritiva, cálculo das métricas de desempenho e realização dos testes estatísticos.


\section{Preparação da base experimental}

\subsection{Contexto}

A preparação da base experimental consistiu na organização dos dados necessários para a aplicação dos modelos de linguagem ao problema de detecção de code smells. 
Para isso, foi utilizado o dataset MLCQ, apresentado por \citeonline{MadeyskiLewowski2020MLCQ}, que reúne amostras de código Java previamente 
classificadas quanto à presença ou ausência de diferentes tipos de code smells. Neste trabalho, foram considerados os smells Blob/God Class, 
Data Class, Feature Envy e Long Method, contemplando tanto elementos de classe quanto elementos de método.

A escolha do MLCQ como fonte de dados experimentais fundamenta-se em evidências 
identificadas na revisão estruturada da literatura. Diversos estudos recentes têm utilizado o 
MLCQ em experimentos envolvendo detecção automática de \textit{code smells} 
\cite{Kovacevic2022AutomaticDetection,Ho2025EnseSmells,Sandouka2025CrossLanguage}. 
Essa adoção decorre do fato de que o dataset foi construído com a participação de desenvolvedores profissionais, 
cujas anotações foram validadas por múltiplos revisores, assegurando confiabilidade e reduzindo vieses individuais, 
além de abranger uma grande diversidade de sistemas reais.

\textbf{Sujeito do experimento:} Os sujeitos do experimento são os arquivos de código-fonte presentes nos projetos que compõem o \textit{dataset MLCQ}.

\textbf{População:} a população é composta pelos 792 projetos Java provenientes do
dataset MLCQ.

\textbf{Amostra:} a amostra é formada pelos aproximadamente 15.000 trechos de código-fonte individualmente anotados
no MLCQ. Cada trecho corresponde a uma instância candidata de \textit{code smell}, definida pelos
metadados do dataset, incluindo: \textit{repository}, \textit{commit hash}, \textit{path}, \textit{start line},
\textit{end line}, \textit{smell type}, \textit{severity} e \textit{type}.

Para a execução do experimento, cada trecho de código será reconstruído a partir dos metadados fornecidos pelo dataset
(\textit{repository}, \textit{commit hash}, \textit{path}, \textit{start line} e \textit{end line}),
garantindo a fidelidade ao código original existente no momento da anotação. As abordagens de detecção serão então aplicadas sobre o mesmo conjunto de trechos de código-fonte,
permitindo comparações entre abordagens distintas de detecção de \textit{code smells}.

Como os tratamentos envolvem modelos fechados (\textit{LLMs}), preserva-se a versão do modelo,
as configurações utilizadas e os prompts empregados, de forma a garantir reprodutibilidade.
Além disso, o experimento não envolve sujeitos humanos, ocorrendo exclusivamente no
ambiente técnico composto pelo dataset, repositórios GitHub e ferramentas automatizadas.


\section{Preparação dos Instrumentos}

A instrumentação se refere às ferramentas e artefatos que serão utilizados para a
execução do experimento. Serão utilizados os seguintes instrumentos para apoiar a execução do experimento:

\begin{itemize}
    \item \textbf{Dataset MLCQ}\footnote{Disponível em: \href{https://doi.org/10.5281/zenodo.3590101}{https://doi.org/10.5281/zenodo.3590101}. Acesso em: 6 dez. 2025.}: o MLCQ é um conjunto 
    de dados de código Java anotado manualmente por desenvolvedores quanto à presença, tipo e severidade de \textit{code smells}. 
    No experimento, o MLCQ será utilizado como fonte dos metadados de cada instância de \textit{code smell}. Cada linha do dataset 
    corresponde a uma anotação individual, identificada pelo campo \textit{id}, associada a um revisor específico (\textit{reviewer\_id}) e 
    a uma amostra de código (\textit{sample\_id}). Os campos \textit{smell} e \textit{severity} registram, respectivamente, o tipo de \textit{code smell} 
    identificado e a severidade atribuída, enquanto \textit{review\_timestamp} armazena a data e hora da anotação. 
    
    O tipo de elemento de código avaliado, por exemplo, função, método ou classe é descrito em \textit{type}, e o nome qualificado desse elemento é registrado em \textit{code\_name}. 
    A origem do artefato é caracterizada pelos campos \textit{repository} (repositório Git de origem), \textit{commit\_hash} (identificador do 
    \textit{commit} em que o trecho foi anotado) e \textit{path} (caminho do arquivo dentro do repositório), enquanto a região de interesse é delimitada 
    por \textit{start\_line} e \textit{end\_line}. O campo \textit{link} fornece a URL para visualização do arquivo e das linhas correspondentes na 
    interface web do GitHub, e \textit{is\_from\_industry\_relevant\_project} indica se o projeto é considerado relevante para a indústria. Em conjunto, 
    esses metadados serão utilizados para localizar o arquivo completo, reconstruir o contexto do trecho anotado e empregar o rótulo humano como 
    referência na avaliação dos tratamentos experimentais.


    \item \textbf{Git}\footnote{Disponível em: \href{https://git-scm.com/}{https://git-scm.com/}. Acesso em: 6 dez. 2025.}: o Git é um sistema de controle de versão distribuído amplamente
    utilizado para gerenciamento de código-fonte. No experimento, o Git será empregado
    para clonar os repositórios necessários, realizar \textit{checkout} em commits
    específicos indicados no MLCQ e assegurar a integridade e consistência dos trechos
    de código reconstruídos.
    
    \item \textbf{GitHub}\footnote{Disponível em: \href{https://github.com/}{https://github.com/}. Acesso em: 6 dez. 2025.}: o GitHub é uma plataforma de hospedagem e versionamento de
    código baseada em Git, amplamente utilizada para projetos open-source. A partir
    das informações fornecidas pelo MLCQ, os repositórios originais serão acessados 
    (via \textit{clone} local ou API do GitHub) para obtenção dos arquivos 
    correspondentes a cada \textit{commit}. Os trechos de código analisados no 
    experimento serão extraídos com base nas linhas inicial e final indicadas no dataset.
    
    \item \textbf{Scripts em Python}: A linguagem Python\footnote{Documentação oficial do Python disponível em: \href{https://www.python.org/}{https://www.python.org/}. Acesso em: 6 dez. 2025.}
    será responsável por automatizar a reconstrução dos trechos a partir do MLCQ e do GitHub, pela execução dos tratamentos 
    (SonarQube e LLMs) e pela comparação das respostas com a rotulação humana.
    
    \item \textbf{SonarQube}\footnote{Disponível em: \href{https://www.sonarqube.org/}{https://www.sonarqube.org/}. Acesso em: 6 dez. 2025.}: ferramenta de Análise Estática configurada com o \textit{rule set} padrão para Java, utilizada como um dos tratamentos de 
    detecção de \textit{code smells}.
    
    \item \textbf{ChatGPT}\footnote{Documentação em: \href{https://platform.openai.com/docs}{https://platform.openai.com/docs}. Acesso em: 6 dez. 2025.}:
    o ChatGPT é um modelo de linguagem de grande porte desenvolvido pela OpenAI,
    baseado na arquitetura \textit{transformer}. Ele é capaz de interpretar instruções
    em linguagem natural e analisar trechos de código, gerando respostas descritivas
    ou classificações. No experimento, o ChatGPT será utilizado com prompts
    padronizados para fornecer uma classificação binária (presença ou ausência de
    \textit{code smell}) para cada trecho de código reconstruído a partir do MLCQ.

    \item \textbf{Gemini}\footnote{Documentação em: \href{https://ai.google.dev/}{https://ai.google.dev/}. Acesso em: 6 dez. 2025.}:
    o Gemini é um modelo de linguagem de grande porte desenvolvido pelo Google,
    também fundamentado na arquitetura \textit{transformer}. Ele é projetado para lidar
    com tarefas complexas envolvendo compreensão, geração e análise de linguagem.
    No experimento, o Gemini será igualmente acessado via prompts padronizados,
    produzindo uma classificação binária sobre a presença ou ausência de 
    \textit{code smell} em cada trecho avaliado.

    \item \textbf{Copilot}\footnote{Documentação em: 
    \href{https://learn.microsoft.com/azure/ai-services/openai/}{https://learn.microsoft.com/azure/ai-services/openai/}. 
    Acesso em: 6 dez. 2025.}:
    refere-se ao modelo de linguagem utilizado 
    pelo Microsoft Copilot e disponibilizado de forma programática por meio da Azure OpenAI API, 
    permitindo sua execução automatizada e reprodutível. Esses modelos, derivados da família GPT, 
    são otimizados para tarefas de programação e análise estrutural de código. 
    No experimento, o modelo será acessado mediante prompts padronizados, fornecendo uma 
    classificação quanto à presença ou ausência de \textit{code smells} para cada trecho de código 
    reconstruído a partir do MLCQ, possibilitando comparação direta com as demais abordagens.


    \item \textbf{DeepSeek}\footnote{Documentação em: \href{https://www.deepseek.com/}{https://www.deepseek.com/}. Acesso em: 6 dez. 2025.}:
    o DeepSeek é um modelo de linguagem de grande porte otimizado para tarefas de
    programação, análise de código e raciocínio técnico. Baseado em uma arquitetura
    avançada de \textit{transformers} customizados, o modelo oferece alta capacidade de
    inspeção estrutural de código e interpretação semântica. No experimento, o DeepSeek 
    será empregado por meio de \textit{prompts} padronizados para realizar a detecção de 
    \textit{code smells}, fornecendo uma classificação comparável à dos demais \textit{LLMs} 
    e permitindo avaliar seu desempenho relativo nas tarefas de analisabilidade.

    
    \item \textbf{Ferramentas de análise estatística}: serão utilizadas planilhas eletrônicas e scripts Python com bibliotecas como 
    \textit{pandas}\footnote{Disponível em: \href{https://pandas.pydata.org/}{https://pandas.pydata.org/}. Acesso em: 6 dez. 2025.} e
    \textit{numpy}\footnote{Disponível em: \href{https://numpy.org/}{https://numpy.org/}. Acesso em: 6 dez. 2025.}
    para cálculo das métricas, construção das distribuições e execução dos testes estatísticos necessários.
\end{itemize}

\subsection{Procedimento de aplicação dos instrumentos}

\begin{enumerate}
    \item \textbf{Leitura da instância no MLCQ.} 
    O \textbf{Dataset MLCQ} é utilizado como fonte de consulta para cada instância anotada de \textit{code smell}, tomando seus metadados como ponto de partida para a análise. A partir dos campos presentes no dataset (como \textit{link} do repositório, \textit{path}, \textit{start line} e \textit{end line}), identificam-se o arquivo em que o trecho está localizado e a região específica do código de interesse.

    \item \textbf{Recuperação do código via Git/GitHub.}
    A partir do campo \textit{link}, acessa-se a URL do repositório no \textbf{GitHub} e utiliza-se o \textbf{Git} para 
    realizar o \textit{clone} (ou atualizar um clone existente) e extrair o arquivo indicado em \textit{path}. Em seguida, o arquivo completo é carregado, e as linhas compreendidas entre \textit{start line} e \textit{end line} são apenas marcadas como referência de localização do \textit{code smell}. Ou seja, o insumo para as ferramentas de detecção é o \textbf{arquivo inteiro} em que o trecho se encontra, e não somente o fragmento delimitado pelas linhas.

    \item \textbf{Validação, registro e randomização dos tratamentos.}
    O arquivo recuperado, bem como a região correspondente ao trecho anotado, são validados. Arquivos ou anotações inconsistentes são descartados. Para os casos válidos, scripts em \textbf{Python} registram o arquivo, seus metadados e a faixa de linhas associada à instância de \textit{code smell} em uma estrutura tabular (\textit{DataFrame} do \textit{pandas}) e randomizam a ordem dos cinco tratamentos (\textbf{SonarQube}, \textbf{ChatGPT}, 
    \textbf{Gemini}, \textbf{Copilot} e \textbf{DeepSeek}) a serem aplicados, em conformidade com o \textit{Completely Randomized Design} adotado.

    \item \textbf{Aplicação do tratamento baseado em análise estática (SonarQube).}
    O arquivo de código é então submetido ao \textbf{SonarQube} e avaliado com o \textit{rule set} padrão para Java. As \textit{issues} 
    classificadas como \textit{code smells} são inspecionadas e, para aquela instância do MLCQ, registra-se uma decisão binária indicando se houve ou não detecção de \textit{code smell} na região do arquivo correspondente às linhas \textit{start line}–\textit{end line}.

    \item \textbf{Aplicação dos tratamentos baseados em LLMs.}
    Em seguida, é construído um \textit{prompt} padronizado contendo instruções em linguagem natural, 
    o conteúdo completo do arquivo Java recuperado e a indicação da região de interesse associada à 
    instância do MLCQ. Esse \textit{prompt} é enviado, na ordem aleatória definida, aos modelos de 
    linguagem avaliados por meio de suas \textit{APIs} oficiais, integradas aos scripts em Python:

    \begin{itemize}
        \item \textbf{ChatGPT} (GPT-5.1);
        \item \textbf{Gemini} (Gemini 3 Pro);
        \item \textbf{Copilot} (GPT-5);
        \item \textbf{DeepSeek} (DeepSeek-V3.2).
    \end{itemize}

    Em todos os casos, serão adotadas configurações controladas de inferência (por exemplo,
    limite de \textit{tokens} e modo de resposta) de forma a reduzir a variabilidade das saídas e favorecer a
    reprodutibilidade. Para cada modelo, a resposta textual é armazenada integralmente e, em etapa
    posterior, processada para gerar uma classificação binária padronizada (presença ou ausência de 
    \textit{code smell} para aquela instância), conforme o formato exigido no \textit{prompt}.


    \item \textbf{Consolidação das saídas por trecho.}
    Para cada instância anotada de \textit{code smell}, os scripts em Python consolidam, em uma única linha de um \textit{DataFrame} (\textit{pandas}), os metadados do 
    \textbf{Dataset MLCQ} (inclusive o rótulo humano) e as decisões binárias produzidas por cada tratamento (\textbf{SonarQube}, \textbf{ChatGPT}, 
    \textbf{Gemini}, \textbf{Copilot} e \textbf{DeepSeek}). A partir dessa consolidação, são derivadas variáveis auxiliares por instância, como acerto, 
    falso positivo, falso negativo e verdadeiro negativo para cada ferramenta.

    \item \textbf{Agregação por projeto e cálculo das variáveis dependentes.}
    Por fim, utilizando as \textbf{ferramentas de análise estatística} (scripts Python com \textit{pandas} e \textit{numpy}, eventualmente apoiados 
    por planilhas eletrônicas), os resultados por instância são agregados por projeto, obtendo-se, para cada tratamento, a 
    \textbf{quantidade de \textit{code smells} identificados por projeto}, variável dependente definida na Seção de Variáveis do Experimento. 
    A partir dessas distribuições, são calculadas as estatísticas descritivas (medianas, variâncias, desvios padrão) e preparados os conjuntos 
    de dados de entrada para o teste de hipóteses (Kruskal--Wallis) descrito anteriormente.
\end{enumerate}



\section{Preparação do ambiente experimental}

A avaliação da validade do experimento segue a classificação discutida e adaptada 
para Engenharia de Software por \citeonline{Wohlin:2012:ESE:2349018}, contemplando quatro categorias principais de validade: conclusão, interna, 
de construto e externa. Nesta subseção, são discutidas as principais ameaças identificadas em cada categoria, bem como as estratégias adotadas para 
mitigá-las no contexto do experimento orientado à tecnologia aqui planejado.

\subsubsection{LLMs}

A validade de construto refere-se ao grau em que o experimento reflete adequadamente os conceitos teóricos que pretende investigar, neste caso, o papel da 
analisabilidade na detecção de code smells e, de forma mais ampla, a contribuição de sistemas baseados em IA para a garantia da qualidade de software. 
Uma primeira ameaça é a \textit{inadequate preoperational explication of constructs}: o construto “capacidade de detecção de code smells” é 
operacionalizado, neste estudo, pela quantidade de code smells identificados por projeto para cada tratamento, medida derivada das comparações entre 
as saídas das ferramentas e os rótulos humanos do MLCQ. Essa escolha enfatiza o volume de detecções corretas mas não captura outros aspectos relevantes, como 
o esforço necessário para interpretar as saídas, a severidade dos problemas encontrados. Tais dimensões são reconhecidas como parte do construto mais amplo de 
analisabilidade, mas não são diretamente abordadas na operacionalização.

\subsubsection{SonarQube}

A validade de construto refere-se ao grau em que o experimento reflete adequadamente os conceitos teóricos que pretende investigar, neste caso, o papel da 
analisabilidade na detecção de code smells e, de forma mais ampla, a contribuição de sistemas baseados em IA para a garantia da qualidade de software. 
Uma primeira ameaça é a \textit{inadequate preoperational explication of constructs}: o construto “capacidade de detecção de code smells” é 
operacionalizado, neste estudo, pela quantidade de code smells identificados por projeto para cada tratamento, medida derivada das comparações entre 
as saídas das ferramentas e os rótulos humanos do MLCQ. Essa escolha enfatiza o volume de detecções corretas mas não captura outros aspectos relevantes, como 
o esforço necessário para interpretar as saídas, a severidade dos problemas encontrados. Tais dimensões são reconhecidas como parte do construto mais amplo de 
analisabilidade, mas não são diretamente abordadas na operacionalização.

\section{Execução do Experimento}

A validade interna está relacionada à existência de relações causais entre tratamentos e resultados, isto é, situações em que uma diferença observada entre 
ferramentas não é devida ao tratamento em si, mas a fatores não controlados. No presente estudo, por se tratar de um experimento orientado à tecnologia, 
sem a participação de sujeitos humanos, diversas ameaças clássicas, como maturação, mortalidade de sujeitos, efeitos sociais entre grupos ou seleção de 
participantes são naturalmente mitigadas.

\section{Coleta e armazenamento das respostas}

A validade de conclusão diz respeito à correção das inferências estatísticas sobre a relação entre os tratamentos e os resultados observados. 
Uma ameaça relevante é o \textit{low statistical power}, isto é, a possibilidade de o experimento não detectar diferenças existentes entre os tratamentos. 
Essa ameaça é mitigada, em parte, pelo tamanho da amostra: os resultados são agregados a partir de aproximadamente 15.000 instâncias anotadas no MLCQ, 
distribuídas em mais de 700 projetos Java, o que tende a aumentar o poder estatístico das análises.



